%!TEX root = ../LLMSeminar.tex
% ──────────────────────────────────────────────────────────────
%  Section 10 — Towards Ambitious Science
% ──────────────────────────────────────────────────────────────

\section{Towards ambitious science}\label{sec:ambitious-science}

\subsection{The lack-of-imagination problem}

Spend any amount of time around the people most active in
the agentic-AI space and a strange pattern is unmissable: the
most obvious thing they all build is \emph{tooling for other
people building agentic AI}.  Frameworks, dashboards,
``observability platforms,'' marketplaces of agents.  The
recursion is not accidental.  It is what one does when one
has acquired a powerful tool but has not yet decided what
the tool is \emph{for}.

\begin{keyresult}[The bottleneck is taste, not capability]
  At the present moment of the field, the rate-limiting step
  in deploying these systems usefully is no longer engineering
  power.  It is \emph{the choice of problem}: what is genuinely
  worth doing, and what would be merely incremental?  Almost
  everyone is under-ambitious here, including the people at
  the frontier of the technology.
\end{keyresult}

The reason is simple.  Until recently, what one could attempt
on a one-year timescale with a small group was rigidly
bounded by the cost of specialist labour.  We learned what
problems to consider \emph{tractable} by absorbing this
boundary.  When the boundary moves, our learned sense of the
tractable does not move with it; we keep proposing the same
size of problem because that is what feels reasonable.  But
the constraints have changed.

\subsection{A reference scale: star formation}

Consider a problem that, until recently, no group of fewer
than several hundred physicists would have proposed to take
on as a single project: a self-consistent, full-physics
simulation of a star's formation, from the first
gravitational instability of a cold molecular cloud to the
ignition of nuclear burning.

\begin{intuition}[All of physics, in one problem]
  Star formation is a deliberately maximalist example
  because the physics genuinely does span almost everything
  taught in an undergraduate degree.  In approximate temporal
  order it requires:
  \begin{itemize}
    \item gravitational collapse and self-gravitating
      hydrodynamics;
    \item magnetohydrodynamics of weakly ionised plasmas
      (ambipolar diffusion in the cool phases);
    \item molecular astro-chemistry and dust formation;
    \item optical and infra-red radiative transfer
      (cooling and heating);
    \item rotational and turbulent angular-momentum
      transport;
    \item nuclear physics for the onset of burning;
    \item and, if one is honest, weak-field general
      relativity for the deepest potentials.
  \end{itemize}
  No single human researcher commands all of these to
  publication standard.  Historically the projects that have
  attempted it---FLASH, Athena++, MFM-based codes---have been
  multi-decade, multi-institution efforts.
\end{intuition}

The thesis of this seminar is that the floor on ``what a
small group with operational competency and good agents can
attempt'' has moved closer to that ceiling.  Not all the
way---that is part of why the question is open---but a long
way.

\begin{keyresult}[Vision for this seminar series]
  We aim to identify a problem of \emph{ambitious} scientific
  scope, attempt it as a group with the help of coding
  agents, and document carefully where the agents accelerate
  the work, where they fail, and what new failure modes
  appear.  Whether or not we succeed, the experiment is
  worth running.
\end{keyresult}

\subsection{An anti-example: $P$ vs $\textsf{NP}$}

\begin{audience}[Audience question]
  \emph{What about asking it to prove $P \neq \textsf{NP}$?}

  This is a natural sceptical reflex: pose an
  acknowledged-hard problem and watch the agent fail.  We
  will encounter this kind of failure repeatedly, and the
  failure mode is informative.  Two observations:
  \begin{enumerate}[(1)]
    \item A proof of $P \neq \textsf{NP}$ may simply not fit
      in the model's context window.  This is the
      Seminar~I limitation re-appearing in a new disguise:
      the function $f$ has bounded input length, so very
      long proofs cannot be expressed in a single forward
      pass even if they could be discovered.
    \item ``One-shotting'' such a problem---feeding the
      statement and asking for a proof in a single
      conversation---is the wrong workflow.  The interesting
      programme is decomposition: \emph{how do we break
      down a problem so that an agent can chip away at it
      over many turns, with verification at every step?}
      That decomposition is the technical question of
      Seminar~III.
  \end{enumerate}
\end{audience}

\subsection{Operational competency vs.\ domain expertise}

A genuine objection arises here, and it deserves careful
treatment.

\begin{audience}[Audience objection (Martin)]
  \emph{To trust an LLM's output you need deep domain
  expertise, because the output is plausible-sounding even
  when wrong.  Without expertise you cannot verify, and you
  will spend more time chasing phantoms than you save.}
\end{audience}

This is a serious point.  The corresponding empirical
observation came up in the room: a tutor in operator theory
on Hilbert spaces fed an exercise to the chat model, and
got a slightly wrong outline of a proof, identifiable as
wrong only by someone who already understood the structure
of symmetric hierarchies in the relevant calculus.  A
non-expert, taking the output at face value, would have
spent days on a dead end.

The question, however, is whether ``expert'' and
``non-expert'' is the correct binary.  I claim it is not.

\begin{definition}[Operational competency]\label{def:opcomp}
  We say a worker has \emph{operational competency} in a
  domain when they can:
  \begin{enumerate}[(i)]
    \item read its primary literature without translation;
    \item identify the standard objects and operations;
    \item run small computations and verify each step against
      a known case;
    \item recognise when a result is suspicious.
  \end{enumerate}
  Operational competency falls short of full domain
  expertise---one cannot, for instance, design a research
  programme from first principles---but it is sufficient to
  \emph{verify} much of what an LLM produces in the domain.
\end{definition}

\begin{intuition}[The verifier-prover asymmetry]
  Many problems---in the technical sense, all of $\textsf{NP}$---are
  asymmetric in difficulty: easy to verify a candidate
  solution, hard to find one.  In such cases the worker
  needs only enough domain skill to \emph{check} steps, not
  to \emph{generate} them.  The LLM proposes; the operationally
  competent human verifies; verification is exponentially
  cheaper than generation when the asymmetry holds.
\end{intuition}

\begin{keyresult}[A working hypothesis]
  Operational competency, plus an effective coding agent,
  is enough to drive substantial progress in many domains
  that previously required full domain expertise of every
  participant.  The cases in which it is \emph{not} enough
  are those where verification itself is hard: where a
  wrong answer cannot be cheaply distinguished from a right
  one.
\end{keyresult}

The unhappy converse follows: agents are most dangerous in
exactly the domains where verification is opaque.  This is
where Martin's objection genuinely applies and where the
``half-assing it'' failure mode becomes serious.  An agent
producing slightly-wrong physical intuition in a sub-field
the user does not understand will quietly poison the
project.

\begin{keyresult}[Diagnostic: is verification cheap here?]
  Before launching a serious agent-driven attack on a
  problem, ask: \emph{can each step's output be cheaply
  checked?}  If the answer is yes, operational competency
  suffices and the project is likely to succeed.  If the
  answer is no, real domain expertise must be at the table.
\end{keyresult}

\subsection{What we will try}

The intent for the remainder of the seminar series is to
take this hypothesis seriously: identify a substantial
physics problem, decompose it under expert supervision into
verifiable pieces, and let coding agents do the bulk of the
implementation.  Candidate domains include large-scale
simulation, formalisation of mathematical proofs (the Lean
ecosystem being a natural fit, as verification is mechanical),
and inverse problems with cheap forward models.  Concrete
proposals from participants are welcome.
